\documentclass[10pt]{article}
\usepackage[a4paper,margin=1in]{geometry}
\usepackage{xeCJK}
\setCJKmainfont{FandolSong-Regular.otf}[ItalicFont=FandolKai-Regular.otf]
\setCJKsansfont{FandolHei-Regular.otf}
\setCJKmonofont{FandolFang-Regular.otf}
\usepackage{hyperref}
\usepackage{enumitem}
\setlist{nosep,leftmargin=1.5em}
\hypersetup{colorlinks=true,linkcolor=blue,urlcolor=blue,citecolor=blue}
\XeTeXlinebreaklocale "zh"
\XeTeXlinebreakskip = 0pt plus 1pt
\setlength{\parindent}{2em}
\setlength{\parskip}{0.35em}
\emergencystretch=3em
\sloppy

\title{Genie: Generative Interactive Environments\\[0.5em]\large 中文重构译文}
\author{Jake Bruce 等著 \quad 中文译文备份}
\date{}

\begin{document}
\maketitle

\noindent\textbf{说明：} 本文档由 PDF 抽取文本重构生成，用于在缺少可用 arXiv LaTeX 源码时备份中文译文。章节和段落为自动恢复，图表、公式和双栏排版未按原始论文完整复刻。

\bigskip

2024-2-26

\section*{Genie：生成式交互环境}
Jake Bruce*, Michael Dennis*, Ashley Edwards*, Jack Parker-Holder*, Yuge (Jimmy) Shi*, Edward Hughes, Matthew Lai, Aditi Mavalankar, Richie Steigerwald, Chris Apps, Yusuf Aytar, Sarah Bechtle, Feryal Behbahani, Stephanie Chan, Nicolas Heess, Lucy Gonzalez, Simon Osindero, Sherjil Ozair, Scott Reed, Jingwei Zhang, Konrad Zolna, Jeff Clune, Nando de Freitas, Satinder Singh 和 Tim Rocktäschel

Google DeepMind；University of British Columbia。arXiv:2402.15391v1 [cs.LG]，2024年2月23日。

图1. 从提示生成可交互环境：Genie 能够把文本、合成图像、照片和手绘草图等不同提示转换为可游玩的环境。该能力来自一个仅从互联网视频中无监督学习得到的潜在动作接口。

\section*{摘要}
本文介绍 Genie。论文将其定义为生成式交互环境，并指出这是首个以无监督方式从无标签互联网视频中训练得到的同类系统。模型可以根据文本、合成图像、照片甚至草图作为提示，生成多种可由动作控制的虚拟世界。Genie 具有 110 亿参数，可以视为一种基础世界模型。它由时空视频 tokenizer、自回归动力学模型，以及简单且可扩展的潜在动作模型组成。

尽管训练过程中没有真实动作标签，也没有世界模型文献中常见的领域特定要求，Genie 仍允许用户在生成环境中逐帧采取动作。学得的潜在动作空间还可用于让智能体从未见过的视频中模仿行为，为后续训练通用智能体提供一个可研究的接口。

关键词：Generative AI，基础模型，世界模型，视频模型，开放性

通讯作者：Ashley Edwards（edwardsashley@google.com）和 Jack Parker-Holder（jparkerholder@google.com）。\textasciicircum{} 2024 Google DeepMind。所有权利保留。

\section*{Genie：生成式交互环境}
\section{导言}
过去几年，生成式 AI 快速发展，模型已经能够生成文本、图像等多种内容。受 Transformer 等架构进展、硬件进步以及模型和数据规模扩大的推动，研究者已经能够生成连贯的对话文本，也能够根据文本提示生成高质量图像。近期结果表明，视频生成也会受益于规模化，并可能成为下一阶段的主要研究方向之一。

不过，现有视频生成模型与 ChatGPT 这类语言工具相比，在交互性和参与感上仍有很大差距，也还难以支持持续交互的体验。一个自然的问题是：给定来自互联网的大规模视频数据，模型能否在生成图像或视频之外，进一步生成完整的交互过程？本文提出“生成式交互环境”这一范式，即从单个文本或图像提示生成可交互环境。

Genie 从超过 20 万小时的公开互联网游戏视频中训练得到；尽管训练时没有动作标注或文本标注，它仍可以通过学得的潜在动作空间进行逐帧控制。在 110 亿参数规模下，Genie 呈现出基础模型的若干特征：它可以把未见过的图像作为提示，创建并游玩由模型生成的虚拟世界。

\section{方法}

Genie 是一个只用视频数据训练的生成式交互环境。本节先介绍模型中使用的基础结构，再说明三个核心组件。Genie 的若干组件基于 Vision Transformer（ViT）。由于视频序列可包含多达 $O(10^4)$ 个 token，普通 Transformer 的二次方内存开销会迅速成为瓶颈。因此，作者在所有组件中采用更节省内存的时空 Transformer（ST-Transformer）结构，在模型容量和计算约束之间取得平衡。

与每个 token 都关注所有其他 token 的传统 Transformer 不同，ST-Transformer 由若干时空块组成，每个块交替使用空间注意力层和时间注意力层，并接一个前馈层（FFW）。空间层在每个时间步内的 $H\times W$ 个 token 上做自注意力；时间层则沿时间维度关注同一空间位置上的 token。时间层采用因果掩码。这样，主要计算量来自空间注意力，并随帧数线性增长，而不是二次增长，因此更适合生成具有长期一致动力学的视频。作者还在 ST 块中只保留空间和时间层之后的一个 FFW，省去空间层之后的额外 FFW，以便把计算预算用于扩大模型的其他部分。

\section*{模型组件}

如图3所示，Genie 包含三个关键组件：潜在动作模型（LAM）在相邻帧之间推断潜在动作；视频 tokenizer 将原始视频帧转换为离散 token；动力学模型在给定潜在动作和历史帧 token 的情况下预测下一帧。训练分为两个阶段：先训练视频 tokenizer，并将其用于后续动力学模型；随后联合训练直接从像素推断潜在动作的 LAM 和基于视频 token 的动力学模型。

图3. Genie 模型训练：模型接收 $T$ 帧视频，通过视频 tokenizer 将其转换为离散 token $z$，并用潜在动作模型推断相邻帧之间的潜在动作 $a$。二者随后输入动力学模型，以自回归方式预测后续帧。

\section*{潜在动作模型（LAM）}

为了实现可控视频生成，Genie 将每个未来帧的预测都条件化在上一帧采取的动作上。然而，互联网视频通常没有动作标签，人工标注也代价很高。因此，作者以完全无监督的方式学习潜在动作。

具体来说，编码器接收历史帧 $x_{1:t}$ 和下一帧 $x_{t+1}$，输出对应的连续潜在动作；解码器再接收历史帧和潜在动作，并预测下一帧。训练时，模型使用基于 VQ-VAE 的目标，将动作限制在一个小的离散码本中。码本大小 $|A|$ 表示可能潜在动作的最大数量，作者在实验中取 $|A|=8$。这个小码本一方面便于人类试玩，另一方面也迫使模型把可控变化压缩进少量动作中。由于解码器只能看到历史和潜在动作，潜在动作必须捕捉过去与未来之间最有意义的变化，才能帮助重建未来帧。

这个解码器只在训练 LAM 时提供学习信号。推断时，除 VQ 码本外，LAM 的其余部分会被丢弃，并由用户提供的动作替代。作者同样用 ST-Transformer 构建 LAM；时间层中的因果掩码允许模型一次接收整段视频 $x_{1:T}$，并生成相邻帧之间的全部潜在动作 $a_{1:T-1}$。

图5. 潜在动作模型：LAM 从无标签视频帧中以无监督方式学习动作。

\section*{视频 Tokenizer}

与先前工作类似，Genie 将视频压缩为离散 token，以降低维度并提升视频生成质量。视频 tokenizer 仍基于 VQ-VAE：给定 $T$ 帧视频 $x_{1:T}$，它为每一帧生成离散表示 $z_{1:T}$，其中 $D$ 表示离散潜在空间的大小。整个视频序列使用标准 VQ-VAE 目标进行训练。

与只在空间维度上压缩的 tokenizer 不同，Genie 在编码器和解码器中都使用 ST-Transformer，把时间动力学纳入编码过程，从而改善视频生成质量。由于 ST-Transformer 具有因果结构，每个离散编码 $z_t$ 都包含此前帧 $x_{1:t}$ 的信息。Phenaki 也使用时间感知 tokenizer C-ViViT，但该结构的计算开销随帧数二次增长；相比之下，Genie 基于 ST-Transformer 的 tokenizer（ST-ViViT）主要计算量随帧数线性增长，因此更高效。

图6. 视频 tokenizer：带 ST-Transformer 的 VQ-VAE。

\section*{动力学模型}

动力学模型是一个仅解码器结构的 MaskGIT Transformer。在时间步 $t\in[1,T]$，它接收已经 token 化的视频 $z_{1:t-1}$ 以及停止梯度的潜在动作 $a_{1:t-1}$，并预测下一帧 token $\hat z_t$。由于仍采用因果 ST-Transformer，模型可以把所有 $T-1$ 帧的 token 和潜在动作作为输入，并一次预测全部下一帧 token $\hat z_{2:T}$。训练目标是预测 token 与真实 token $z_{2:T}$ 之间的交叉熵损失。训练时，输入 token $z_{2:T-1}$ 会按伯努利分布随机掩码，掩码率从 $0.5$ 到 $1$ 均匀采样。

许多世界模型会把时间步 $t$ 的动作拼接到对应帧上。作者发现，将潜在动作作为加性嵌入输入 LAM 和动力学模型，能够提升生成结果的可控性。

图7. 动力学模型：模型接收视频 token 和动作嵌入，并预测未来被掩码的视频 token。图8. Genie 推断：提示帧先被 token 化，再与用户选择的潜在动作结合并输入动力学模型进行迭代生成；预测出的帧 token 最后通过 tokenizer 解码回图像空间。

\section*{推断：动作可控的视频生成}

推断时，玩家首先提供一张图像 $x_1$ 作为初始帧，模型用视频编码器将其转换为 $z_1$。随后，玩家选择一个离散潜在动作 $a_1$，即在 $[0, |A|)$ 中选择一个整数。模型将帧 token $z_1$ 与由离散动作索引得到的潜在动作嵌入一起输入动力学模型，预测下一帧 token $z_2$。该过程不断重复，随着用户继续输入动作，模型自回归生成剩余序列，并由 tokenizer 的解码器把 token 转回视频帧。

刚开始与模型交互时，用户并不知道每个潜在动作会如何影响下一帧。不过作者观察到，同一动作在不同输入上的含义较为一致，因此理解潜在动作的映射关系很像学习一只新手柄的按键。

\section{实验结果}

\paragraph{数据集}

作者在公开互联网视频中收集了大规模 2D 平台游戏数据集（下称 Platformers）来训练 Genie。通过筛选与平台游戏相关的关键词，得到约 5500 万个 16 秒、10 FPS、分辨率为 $160\times90$ 的视频片段。最终数据集包含 680 万个 16 秒视频片段，约 3 万小时，规模与其他常用互联网视频数据集处于同一数量级。除非特别说明，实验结果均来自在该数据集上训练的 110 亿参数模型。

为验证方法的通用性，作者还使用 RT-1 机器人数据集进行训练：该数据集包括约 13 万条机器人演示、一个独立的模拟数据集，以及先前工作中的 20.9 万条真实机器人回合。这里不使用任何动作标签，只把它们当作视频处理。下文将该数据集简称为 Robotics。

\paragraph{指标与训练细节}

作者从视频保真度和可控性两个方面评估 Genie。视频保真度使用 Fréchet Video Distance（FVD）衡量；可控性使用 $\Delta_t$ PSNR，度量在相同历史下改变潜在动作会使生成视频偏离真实未来多少。所有实验报告 $t=4$ 时的 $\Delta_t$ PSNR。

视频 tokenizer 使用 2 亿参数、patch 大小 4、嵌入大小 32、1024 个离散码的码本。潜在动作模型使用 3 亿参数、patch 大小 16、嵌入大小 32、8 个离散码（即 8 个潜在动作）。所有建模组件都使用 16 帧序列，帧率为 10 FPS。训练动力学模型时使用 bfloat16 和 QK norm，并在视频质量评估中采用 25 个 MaskGIT 步骤。更多细节见附录 C。

\subsection{规模化结果}

本节研究模型大小和批量大小的规模化行为。作者将动力学模型从 4000 万参数扩展到 27 亿参数，并在固定 tokenizer 和 LAM 结构下观察训练损失。图9显示，随着计算量增加，最终训练损失持续下降，说明该架构能从更大的模型和批量中受益。最终模型使用 101 亿参数的动力学模型、批量大小 512、125k 步训练，并使用 256 个 TPU v5p；加上 tokenizer 和动作模型后，总参数量约为 107 亿，文中称为 110 亿参数 Genie 模型。

图9. 规模化结果：左侧为不同模型大小的训练曲线；中间为每个模型大小在最近 300 次更新上的平均最终训练损失；右侧为 23 亿参数模型在不同批量大小下的最终训练损失。

\subsection{定性结果}

作者展示了两个模型的定性结果：一个是在 Platformers 数据集上训练的 110 亿参数 Genie 模型，另一个是在 Robotics 数据集上训练的较小模型。整体上，Genie 能在不同域中生成质量较高且可控的视频。Platformers 模型的定性评估只使用分布外图像提示，包括文本到图像模型生成的图像、手绘草图和真实照片。即使这些提示在视觉上明显不同于训练数据，模型仍能生成类似游戏的交互行为，说明大规模视频训练提高了模型对提示变化的适应性。

图10. 从图像提示开始游玩：Genie 可以接收文本到图像模型生成的图像、手绘草图或真实照片作为提示。图中展示提示帧，以及连续四次采取同一潜在动作后的帧。尽管部分图像与训练集差异较大，角色仍呈现清晰移动。

实验中还观察到一些由数据驱动产生的行为。例如，在图12中，给定由 Imagen2 生成的图像后，潜在动作会让前景、中景和背景以不同速度移动，模拟平台游戏中常见的视差效果。作者还训练了一个 25 亿参数的 Robotics 模型，使用与 Platformers 上最优设置相同的超参数，在测试集上达到 82.7 的 FVD。如图13所示，该模型在没有文本和动作标签的情况下，从视频中学到了清晰且一致的机器人动作；除机械臂控制外，模型还能够模拟物体交互和变形，例如图11中芯片袋的物理变化。这说明 Genie 的方法可能为机器人学提供一种基于大规模视频数据的低层级可控世界模型。

\subsection{训练智能体}

作者将 Genie 定位为训练通用智能体时可使用的一类基础世界模型。图14显示，给定未见过的 RL 环境初始帧，Genie 已经能够生成多样轨迹。为进一步验证从互联网视频学得的潜在动作是否可用于模仿未见行为，作者冻结 LAM，用它为目标环境中的专家视频序列标注离散潜在动作；随后训练一个策略，根据观测预测专家采取各个潜在动作的可能性。最后，再用一小批包含真实专家动作的数据，将潜在动作映射到真实动作，细节见附录 E。

实验在程序生成的 2D 平台环境 CoinRun 中进行，分别考虑简单和困难设置。上界是可以访问专家动作的 oracle 行为克隆模型，下界是随机智能体。LAM 策略只需约 200 个专家样本进行适配，就能达到 oracle 行为克隆的水平，尽管模型此前几乎不可能见过 CoinRun。这说明学得的潜在动作具有跨环境的一致性和可迁移意义，因为从潜在动作到真实动作的映射并不包含当前观测信息。

\subsection{消融研究}

在设计潜在动作模型时，作者比较了两类输入：最终采用的原始图像（pixel-input），以及把图像先 token 化后再输入 LAM 的方案（token-input）。虽然 token-input 在 Platformers 上获得略低的 FVD，但在 Robotics 上没有保持优势；从控制指标看，token-input 在两个环境中的可控性都更差，$\Delta_t$ PSNR 更低。这表明 token 化过程可能丢失了部分关于视频动力学和运动的信息，因此让 LAM 直接接收原始视频更有利。

作者还比较了三种 tokenizer 架构：只建模空间的 ViT、时空 ST-ViViT，以及时空 C-ViViT。所有 tokenizer 使用相近参数量、patch 大小 10、批量大小 128 和序列长度 16，然后训练相同的动力学模型和潜在动作模型。结果显示，ST-ViViT 在视频生成质量（FVD）和可控性（$\Delta_t$ PSNR）之间取得了更好的平衡。C-ViViT 使用完整时空注意力，在相同参数量下内存消耗显著更高，但性能并未提升，并表现出过拟合倾向，需要更强正则化。

\section{相关工作}

\paragraph{世界模型}
生成式交互环境可以视为世界模型的一类。传统世界模型通常在动作条件下预测下一帧，可用于在训练智能体时提供想象经验，减少对真实环境交互的依赖。不过，这类模型本身往往需要从环境中直接获得带动作条件的数据。与之不同，Genie 只从公开互联网视频中无监督学习世界模型。近期 GAIA-1 和 UniSim 也强调世界模型规模化，分别面向自动驾驶和机器人操作，但它们需要文本或动作标签；Genie 的重点是从纯视频数据中学习。

\paragraph{视频模型}
Genie 与视频生成模型密切相关。视频模型通常以初始帧或文本为条件，预测视频中的剩余帧。Genie 的动力学模型与 Phenaki、TECO、MaskViT 等基于 Transformer 的视频模型最接近，因为它同样在 token 化图像上使用 MaskGIT 和 ST-Transformer。不同之处在于，Genie 追求更具智能体属性的目标：显式从数据中学习潜在动作空间，使用户或智能体能够通过潜在动作条件预测来“玩”这个模型。

\paragraph{可游玩视频生成}
Genie 扩展了可游玩视频生成（PVG）的设定。PVG 使用潜在动作控制直接从视频学习得到的世界模型，但通常处理特定领域中的静态样例，而不是通过提示生成全新的环境。Genie 为了扩展到更通用的设定，对架构做了实质调整，减少特定归纳偏置，换取更一般的建模能力。

\paragraph{环境生成}
本文也与程序内容生成（PCG）相关。机器学习已经被用于生成游戏关卡，近期也有工作利用语言模型直接编写游戏代码。语言模型本身也可以被看作交互环境，只是缺少视觉部分。Genie 的不同之处在于，它直接从像素学习和生成关卡式环境，因此可以利用互联网视频数据中的多样性。

\paragraph{用潜在动作训练智能体}
已有工作使用潜在动作进行观察模仿、规划和 RL 智能体预训练。这些方法的目标与 Genie 的潜在动作模型相似，但尚未在同等规模上应用。VPT 使用由人类标注动作数据训练出的逆动力学模型，为互联网规模视频打上动作标签，再用这些标签训练策略。Genie 则表明，可以直接利用从互联网视频中学得的潜在动作，在任意环境中推断策略，避免依赖昂贵且可能难以泛化的真实动作标签。

\section{结论与未来工作}

本文提出 Genie，一种新的生成式 AI 形式，使用户能够构想、创建并进入生成的世界，类似于进入人工设计的模拟环境。即使只使用视频数据训练，Genie 仍能根据提示生成多样、交互式、可控的环境。

模型仍有明确改进空间。Genie 继承了自回归 Transformer 的一些弱点，可能幻觉出不真实的未来。虽然时空表示已经取得进展，但当前模型只保留 16 帧记忆，因此在更长时间尺度上保持环境一致性仍有难度。此外，Genie 当前约以 1 FPS 运行，要达到高效交互帧率还需要后续改进。

尽管如此，Genie 为未来研究打开了很多方向。由于方法具有通用性，模型可以在更大比例的互联网视频上训练，以模拟更丰富、真实或想象中的环境。本文只初步讨论了使用 Genie 训练智能体的可能性；考虑到缺乏丰富多样环境是 RL 的关键限制之一，这一路线可能为构建更通用的智能体提供新途径。

\section*{更广泛影响}

\paragraph{社会影响}
Genie 可能让更多人生成自己的类游戏体验。对希望用新方式表达创造力的人来说，这可能是积极的，例如孩子可以设计并进入自己想象的世界。作者同时指出，随着技术进步，必须认真探索如何让这种能力用于扩展人类已有的游戏创作和创意过程，并让相关行业以负责的方式使用 Genie 开发下一代可游玩世界。

\paragraph{训练数据与权重}
作者选择不随论文或网站发布训练后的模型检查点、训练数据集或训练数据样例。他们希望先与研究社区和游戏社区进一步交流，确保未来任何发布都尊重安全和责任边界。

\paragraph{可复现性}
论文指出，对计算资源较少的研究者来说，复现主要结果会比较困难。为缓解这一问题，附录 F 描述了一个较小规模、可完全复现的例子，可在单个中端 TPU 或 GPU 上运行。由于许多设计选择可以在两种规模之间迁移，作者认为这能帮助社区继续研究架构改进和新的研究方向。

\section*{致谢}

作者感谢 Mateusz Malinowski、Philip Ball 和 Louis Kirsch 审阅论文草稿；感谢 Cassidy Hardin、David Bridson、Eric Lau、Lars Lowe Sjoesund、Lucas Smaira 和 Bernardo Avila Pires 帮助构建 Platformers 数据集；感谢 Ruben Villegas 就视频模型训练与评估提供讨论；感谢 Adrian Bolton、Rushil Mistry、Hannah Openshaw、Zoubin Ghahramani、Raia Hadsell、Koray Kavukcuoglu、Daan Wierstra、Doina Precup 和 Ed Hirst 提供战略建议与指导。作者使用了 DeepMind JAX 生态系统，并特别感谢 Andy Brock 构建训练模型所用的内部框架，以及 Arthur Brussee 提供早期界面，使研究者能够“游玩”模型。最后，作者感谢 Seneca 和 Caspian Clune 提供创意草图。

\section*{作者贡献}

作者按姓氏字母顺序列出。通讯请联系 Ashley Edwards（edwardsashley@google.com）和 Jack Parker-Holder（jparkerholder@google.com）。

\paragraph{核心贡献者}
\begin{itemize}
  \item Jake Bruce：项目领导，视频 tokenizer 研究，动作模型研究，动力学模型研究，规模化，模型演示，基础设施。
  \item Michael Dennis：动力学模型研究，规模化，指标，模型演示，基础设施。
  \item Ashley Edwards：Genie 概念，项目领导，动作模型研究，智能体训练，模型演示。
  \item Edward Hughes：动力学模型研究，基础设施。
  \item Matthew Lai：数据集整理，基础设施。
  \item Aditi Mavalankar：动作模型研究，指标，智能体训练。
  \item Jack Parker-Holder：Genie 概念，项目领导，动力学模型研究，规模化，数据集整理。
  \item Yuge (Jimmy) Shi：视频 tokenizer 研究，动力学模型研究，数据集整理，指标。
  \item Richie Steigerwald：数据集整理，指标。
\end{itemize}

\paragraph{部分贡献者、顾问与发起人}
\begin{itemize}
  \item Chris Apps：项目管理。
  \item Yusuf Aytar、Sarah Bechtle、Stephanie Chan、Simon Osindero、Sherjil Ozair、Scott Reed、Jingwei Zhang：技术建议。
  \item Feryal Behbahani、Jeff Clune、Nicolas Heess、Nando de Freitas、Satinder Singh：战略建议。
  \item Lucy Gonzalez：项目管理。
  \item Konrad Zolna：规模化与技术建议。
  \item Tim Rocktäschel：Genie 概念与项目领导。
\end{itemize}

\section*{Genie：生成式交互环境}
\section*{参考资料}
I. Babuschkin, K. Baumli, A. Bell, S. Bhupatiraju, J. Bruce, P. Buchlovsky, D. Budden, T. Cai, A. Clark, I. Danihelka, et al. The deepmind jax ecosystem, 2020. URL http://github. com/deepmind, 2010.

M. Bain, A. Nagrani, G. Varol, and A. Zisserman. Frozen in time: A joint video and image encoder for end-to-end retrieval. In 2021 IEEE/CVF International Conference on Computer Vision (ICCV), pages 1708–1718, Los Alamitos, CA, USA, oct 2021. IEEE Computer Society. doi: 10.1109/ICCV48922. 2021.00175.

B. Baker, I. Akkaya, P. Zhokov, J. Huizinga, J. Tang, A. Ecoffet, B. Houghton, R. Sampedro, and J. Clune. Video pretraining (vpt): Learning to act by watching unlabeled online videos. Advances in Neural Information Processing Systems, 35:24639–24654, 2022.

C. Bamford and S. M. Lucas. Neural game engine: Accurate learning ofgeneralizable forward models from pixels. In Conference on Games, 2020.

J. Bauer, K. Baumli, F. Behbahani, A. Bhoopchand, N. Bradley-Schmieg, M. Chang, N. Clay, A. Collister, V. Dasagi, L. Gonzalez, K. Gregor, E. Hughes, S. Kashem, M. Loks-Thompson, H. Openshaw, J. ParkerHolder, S. Pathak, N. Perez-Nieves, N. Rakicevic, T. Rocktäschel, Y. Schroecker, S. Singh, J. Sygnowski, K. Tuyls, S. York, A. Zacherl, and L. M. Zhang. Human-timescale adaptation in an open-ended task space. In A. Krause, E. Brunskill, K. Cho, B. Engelhardt, S. Sabato, and J. Scarlett, editors, Proceedings of the 40th International Conference on Machine Learning, volume 202 of Proceedings of Machine Learning Research, pages 1887–1935. PMLR, 23–29 Jul 2023.

A. Blattmann, T. Dockhorn, S. Kulal, D. Mendelevitch, M. Kilian, D. Lorenz, Y. Levi, Z. English, V. Voleti, A. Letts, V. Jampani, and R. Rombach. Stable video diffusion: Scaling latent video diffusion models to large datasets, 2023a.

A. Blattmann, R. Rombach, H. Ling, T. Dockhorn, S. W. Kim, S. Fidler, and K. Kreis. Align your latents: High-resolution video synthesis with latent diffusion models. 2023 IEEE/CVF Conference on Computer Vision and Pattern Recognition (CVPR), pages 22563–22575, 2023b.

A. Brohan, N. Brown, J. Carbajal, Y. Chebotar, J. Dabis, C. Finn, K. Gopalakrishnan, K. Hausman, A. Herzog, J. Hsu, J. Ibarz, B. Ichter, A. Irpan, T. Jackson, S. Jesmonth, N. J. Joshi, R. Julian, D. Kalashnikov, Y. Kuang, I. Leal, K.-H. Lee, S. Levine, Y. Lu, U. Malla, D. Manjunath, I. Mordatch, O. Nachum, C. Parada, J. Peralta, E. Perez, K. Pertsch, J. Quiambao, K. Rao, M. Ryoo, G. Salazar, P. Sanketi, K. Sayed, J. Singh, S. Sontakke, A. Stone, C. Tan, H. Tran, V. Vanhoucke, S. Vega, Q. Vuong, F. Xia, T. Xiao, P. Xu, S. Xu, T. Yu, and B. Zitkovich. Rt-1: Robotics transformer for real-world control at scale. In Robotics: Science and Systems, 2023.

T. Brooks, B. Peebles, C. Homes, W. DePue, Y. Guo, L. Jing, D. Schnurr, J. Taylor, T. Luhman, E. Luhman, C. W. Y. Ng, R. Wang, and A. Ramesh. Video generation models as world simulators. 2024. URL https://openai.com/research/video-generation-models-as-world-simulators. T. Brown, B. Mann, N. Ryder, M. Subbiah, J. D. Kaplan, P. Dhariwal, A. Neelakantan, P. Shyam, G. Sastry, A. Askell, et al. Language models are few-shot learners. Advances in neural information processing systems, 33:1877–1901, 2020.

H. Chang, H. Zhang, L. Jiang, C. Liu, and W. T. Freeman. Maskgit: Masked generative image transformer. In Proceedings of the IEEE/CVF Conference on Computer Vision and Pattern Recognition (CVPR), pages 11315–11325, June 2022.

\section*{Genie：生成式交互环境}
S. Chiappa, S. Racaniere, D. Wierstra, and S. Mohamed. Recurrent environment simulators. In International Conference on Learning Representations, 2017.

A. Clark, J. Donahue, and K. Simonyan. Efficient video generation on complex datasets. CoRR, abs/1907.06571, 2019. URL http://arxiv.org/abs/1907.06571.

J. Clune. Ai-gas: Ai-generating algorithms, an alternate paradigm for producing general artificial intelligence. arXiv preprint arXiv:1905.10985, 2019.

K. Cobbe, C. Hesse, J. Hilton, and J. Schulman. Leveraging procedural generation to benchmark reinforcement learning. In Proceedings of the 37th International Conference on Machine Learning, pages 2048–2056, 2020.

M. Dehghani, J. Djolonga, B. Mustafa, P. Padlewski, J. Heek, J. Gilmer, A. P. Steiner, M. Caron, R. Geirhos, I. Alabdulmohsin, R. Jenatton, L. Beyer, M. Tschannen, A. Arnab, X. Wang, C. Riquelme Ruiz, M. Minderer, J. Puigcerver, U. Evci, M. Kumar, S. V. Steenkiste, G. F. Elsayed, A. Mahendran, F. Yu, A. Oliver, F. Huot, J. Bastings, M. Collier, A. A. Gritsenko, V. Birodkar, C. N. Vasconcelos, Y. Tay, T. Mensink, A. Kolesnikov, F. Pavetic, D. Tran, T. Kipf, M. Lucic, X. Zhai, D. Keysers, J. J. Harmsen, and N. Houlsby. Scaling vision transformers to 22 billion parameters. In A. Krause, E. Brunskill, K. Cho, B. Engelhardt, S. Sabato, and J. Scarlett, editors, Proceedings of the 40th International Conference on Machine Learning, volume 202 of Proceedings of Machine Learning Research, pages 7480–7512. PMLR, 23–29 Jul 2023.

A. Dosovitskiy, L. Beyer, A. Kolesnikov, D. Weissenborn, X. Zhai, T. Unterthiner, M. Dehghani, M. Minderer, G. Heigold, S. Gelly, J. Uszkoreit, and N. Houlsby. An image is worth 16x16 words: Transformers for image recognition at scale. In International Conference on Learning Representations, 2021. URL https://openreview.net/forum?id=YicbFdNTTy.

A. Edwards, H. Sahni, Y. Schroecker, and C. Isbell. Imitating latent policies from observation. In International conference on machine learning, pages 1755–1763. PMLR, 2019.

S. M. A. Eslami, D. J. Rezende, F. Besse, F. Viola, A. S. Morcos, M. Garnelo, A. Ruderman, A. A. Rusu, I. Danihelka, K. Gregor, D. P. Reichert, L. Buesing, T. Weber, O. Vinyals, D. Rosenbaum, N. Rabinowitz, H. King, C. Hillier, M. Botvinick, D. Wierstra, K. Kavukcuoglu, and D. Hassabis. Neural scene representation and rendering. Science, 360(6394):1204–1210, 2018. doi: 10.1126/ science.aar6170.

P. Esser, J. Chiu, P. Atighehchian, J. Granskog, and A. Germanidis. Structure and content-guided video synthesis with diffusion models. In 2023 IEEE/CVF International Conference on Computer Vision (ICCV), 2023.

C. Finn, I. Goodfellow, and S. Levine. Unsupervised learning for physical interaction through video prediction. In Proceedings of the 30th International Conference on Neural Information Processing Systems, NIPS’16, page 64–72, Red Hook, NY, USA, 2016. Curran Associates Inc. ISBN 9781510838819.

A. Gupta, S. Tian, Y. Zhang, J. Wu, R. Martín-Martín, and L. Fei-Fei. Maskvit: Masked visual pretraining for video prediction. In The Eleventh International Conference on Learning Representations, 2023.

D. Ha and J. Schmidhuber. Recurrent world models facilitate policy evolution. In Proceedings of the 32Nd International Conference on Neural Information Processing Systems, NeurIPS’18, pages 2455–2467, 2018.

\section*{Genie：生成式交互环境}
D. Hafner, T. Lillicrap, J. Ba, and M. Norouzi. Dream to control: Learning behaviors by latent imagination. In International Conference on Learning Representations, 2020.

D. Hafner, T. P. Lillicrap, M. Norouzi, and J. Ba. Mastering atari with discrete world models. In International Conference on Learning Representations, 2021.

K. He, X. Zhang, S. Ren, and J. Sun. Deep residual learning for image recognition. In 2016 IEEE Conference on Computer Vision and Pattern Recognition (CVPR), pages 770–778, 2016. doi: 10.1109/ CVPR.2016.90.

A. Henry, P. R. Dachapally, S. S. Pawar, and Y. Chen. Query-key normalization for transformers. In Findings of the Association for Computational Linguistics: EMNLP 2020, pages 4246–4253, Online, Nov. 2020. Association for Computational Linguistics. doi: 10.18653/v1/2020.findings-emnlp.379.

J. Ho, W. Chan, C. Saharia, J. Whang, R. Gao, A. Gritsenko, D. P. Kingma, B. Poole, M. Norouzi, D. J. Fleet, and T. Salimans. Imagen video: High definition video generation with diffusion models, 2022a.

J. Ho, T. Salimans, A. Gritsenko, W. Chan, M. Norouzi, and D. J. Fleet. Video diffusion models. In S. Koyejo, S. Mohamed, A. Agarwal, D. Belgrave, K. Cho, and A. Oh, editors, Advances in Neural Information Processing Systems, volume 35, pages 8633–8646. Curran Associates, Inc., 2022b.

W. Hong, M. Ding, W. Zheng, X. Liu, and J. Tang. Cogvideo: Large-scale pretraining for text-to-video generation via transformers. arXiv preprint arXiv:2205.15868, 2022.

W. Hong, M. Ding, W. Zheng, X. Liu, and J. Tang. Cogvideo: Large-scale pretraining for text-to-video generation via transformers. In The Eleventh International Conference on Learning Representations, 2023. URL https://openreview.net/forum?id=rB6TpjAuSRy.

T. Höppe, A. Mehrjou, S. Bauer, D. Nielsen, and A. Dittadi. Diffusion models for video prediction and infilling. Transactions on Machine Learning Research, 2022. ISSN 2835-8856.

A. Hu, L. Russell, H. Yeo, Z. Murez, G. Fedoseev, A. Kendall, J. Shotton, and G. Corrado. Gaia-1: A generative world model for autonomous driving, 2023.

J. Huang, Y. Jin, K. M. Yi, and L. Sigal. Layered controllable video generation. In Computer Vision – ECCV 2022: 17th European Conference, Tel Aviv, Israel, October 23–27, 2022, Proceedings, Part XVI, page 546–564, Berlin, Heidelberg, 2022. Springer-Verlag. ISBN 978-3-031-19786-4.

N. P. Jouppi, D. H. Yoon, G. Kurian, S. Li, N. Patil, J. Laudon, C. Young, and D. Patterson. A domainspecific supercomputer for training deep neural networks. Communications of the ACM, 63(7): 67–78, 2020.

D. Kalashnikov, A. Irpan, P. Pastor, J. Ibarz, A. Herzog, E. Jang, D. Quillen, E. Holly, M. Kalakrishnan, V. Vanhoucke, et al. Qt-opt: Scalable deep reinforcement learning for vision-based robotic manipulation. arXiv preprint arXiv:1806.10293, 2018.

N. Kalchbrenner, A. van den Oord, K. Simonyan, I. Danihelka, O. Vinyals, A. Graves, and K. Kavukcuoglu. Video pixel networks. In D. Precup and Y. W. Teh, editors, Proceedings of the 34th International Conference on Machine Learning, volume 70 of Proceedings of Machine Learning Research, pages 1771–1779. PMLR, 06–11 Aug 2017. URL https://proceedings.mlr.press/ v70/kalchbrenner17a.html.

\section*{Genie：生成式交互环境}
S. Kapturowski, G. Ostrovski, J. Quan, R. Munos, and W. Dabney. Recurrent experience replay in distributed reinforcement learning. In International conference on learning representations, 2018.

S. W. Kim, Y. Zhou, J. Philion, A. Torralba, and S. Fidler. Learning to simulate dynamic environments with gamegan. In Proceedings of the IEEE/CVF Conference on Computer Vision and Pattern Recognition (CVPR), June 2020.

S. W. Kim, J. Philion, A. Torralba, and S. Fidler. Drivegan: Towards a controllable high-quality neural simulation. In Proceedings of the IEEE/CVF Conference on Computer Vision and Pattern Recognition (CVPR), pages 5820–5829, June 2021.

G. Le Moing, J. Ponce, and C. Schmid. Ccvs: Context-aware controllable video synthesis. In M. Ranzato, A. Beygelzimer, Y. Dauphin, P. Liang, and J. W. Vaughan, editors, Advances in Neural Information Processing Systems, volume 34, pages 14042–14055. Curran Associates, Inc., 2021.

W. Lotter, G. Kreiman, and D. Cox. Deep predictive coding networks for video prediction and unsupervised learning. In International Conference on Learning Representations, 2017. URL https: //openreview.net/forum?id=B1ewdt9xe. P. Luc, A. Clark, S. Dieleman, D. de Las Casas, Y. Doron, A. Cassirer, and K. Simonyan. Transformationbased adversarial video prediction on large-scale data. CoRR, abs/2003.04035, 2020.

W. Menapace, S. Lathuilière, S. Tulyakov, A. Siarohin, and E. Ricci. Playable video generation. In IEEE Conference on Computer Vision and Pattern Recognition, CVPR 2021, virtual, June 19-25, 2021, pages 10061–10070. Computer Vision Foundation / IEEE, 2021.

W. Menapace, S. Lathuilière, A. Siarohin, C. Theobalt, S. Tulyakov, V. Golyanik, and E. Ricci. Playable environments: Video manipulation in space and time. In Proceedings of the IEEE/CVF Conference on Computer Vision and Pattern Recognition, 2022.

V. Micheli, E. Alonso, and F. Fleuret. Transformers are sample-efficient world models. In The Eleventh International Conference on Learning Representations, 2023.

M. S. Nunes, A. Dehban, P. Moreno, and J. Santos-Victor. Action-conditioned benchmarking of robotic video prediction models: a comparative study. In 2020 IEEE International Conference on Robotics and Automation (ICRA), pages 8316–8322, 2020. doi: 10.1109/ICRA40945.2020.9196839.

J. Oh, X. Guo, H. Lee, R. Lewis, and S. Singh. Action-conditional video prediction using deep networks in atari games. In Proceedings of the 28th International Conference on Neural Information Processing Systems - Volume 2, NIPS’15, page 2863–2871, Cambridge, MA, USA, 2015. MIT Press.

Open Ended Learning Team, A. Stooke, A. Mahajan, C. Barros, C. Deck, J. Bauer, J. Sygnowski, M. Trebacz, M. Jaderberg, M. Mathieu, N. McAleese, N. Bradley-Schmieg, N. Wong, N. Porcel, R. Raileanu, S. Hughes-Fitt, V. Dalibard, and W. M. Czarnecki. Open-ended learning leads to generally capable agents. CoRR, abs/2107.12808, 2021.

M. Oquab, T. Darcet, T. Moutakanni, H. Vo, M. Szafraniec, V. Khalidov, P. Fernandez, D. Haziza, F. Massa, A. El-Nouby, et al. Dinov2: Learning robust visual features without supervision. arXiv preprint arXiv:2304.07193, 2023.

M. Pan, X. Zhu, Y. Wang, and X. Yang. Iso-dream: Isolating and leveraging noncontrollable visual dynamics in world models. In S. Koyejo, S. Mohamed, A. Agarwal, D. Belgrave, K. Cho, and A. Oh, editors, Advances in Neural Information Processing Systems, volume 35, pages 23178–23191. Curran Associates, Inc., 2022.

\section*{Genie：生成式交互环境}
A. Radford, K. Narasimhan, T. Salimans, and I. Sutskever. Improving language understanding by generative pre-training. 2018.

A. Radford, J. Wu, R. Child, D. Luan, D. Amodei, I. Sutskever, et al. Language models are unsupervised multitask learners. OpenAI blog, 1(8):9, 2019.

S. Rajbhandari, J. Rasley, O. Ruwase, and Y. He. Zero: Memory optimizations toward training trillion parameter models. In SC20: International Conference for High Performance Computing, Networking, Storage and Analysis, pages 1–16. IEEE, 2020.

A. Ramesh, M. Pavlov, G. Goh, S. Gray, C. Voss, A. Radford, M. Chen, and I. Sutskever. Zero-shot text-to-image generation. In M. Meila and T. Zhang, editors, Proceedings of the 38th International Conference on Machine Learning, volume 139 of Proceedings of Machine Learning Research, pages 8821–8831. PMLR, 18–24 Jul 2021.

A. Ramesh, P. Dhariwal, A. Nichol, C. Chu, and M. Chen. Hierarchical text-conditional image generation with clip latents, 2022.

S. Reed, K. Zolna, E. Parisotto, S. G. Colmenarejo, A. Novikov, G. Barth-maron, M. Giménez, Y. Sulsky, J. Kay, J. T. Springenberg, T. Eccles, J. Bruce, A. Razavi, A. Edwards, N. Heess, Y. Chen, R. Hadsell, O. Vinyals, M. Bordbar, and N. de Freitas. A generalist agent. Transactions on Machine Learning Research, 2022. ISSN 2835-8856. Featured Certification, Outstanding Certification.

S. Risi and J. Togelius. Increasing generality in machine learning through procedural content generation. Nature Machine Intelligence, 2, 08 2020a. doi: 10.1038/s42256-020-0208-z.

S. Risi and J. Togelius. Procedural content generation: From automatically generating game levels to increasing generality in machine learning. Nature, 2020b.

J. Robine, M. Höftmann, T. Uelwer, and S. Harmeling. Transformer-based world models are happy with 100k interactions. In The Eleventh International Conference on Learning Representations, 2023.

R. Rombach, A. Blattmann, D. Lorenz, P. Esser, and B. Ommer. High-resolution image synthesis with latent diffusion models. In Proceedings of the IEEE/CVF Conference on Computer Vision and Pattern Recognition (CVPR), pages 10684–10695, June 2022.

O. Rybkin*, K. Pertsch*, K. G. Derpanis, K. Daniilidis, and A. Jaegle. Learning what you can do before doing anything. In International Conference on Learning Representations, 2019.

C. Saharia, W. Chan, S. Saxena, L. Li, J. Whang, E. Denton, S. K. S. Ghasemipour, R. Gontijo-Lopes, B. K. Ayan, T. Salimans, J. Ho, D. J. Fleet, and M. Norouzi. Photorealistic text-to-image diffusion models with deep language understanding. In A. H. Oh, A. Agarwal, D. Belgrave, and K. Cho, editors, Advances in Neural Information Processing Systems, 2022.

D. Schmidt and M. Jiang. Learning to act without actions. In The Twelfth International Conference on Learning Representations, 2024.

M. Shoeybi, M. Patwary, R. Puri, P. LeGresley, J. Casper, and B. Catanzaro. Megatron-lm: Training multi-billion parameter language models using model parallelism. CoRR, abs/1909.08053, 2019. URL http://arxiv.org/abs/1909.08053.

U. Singer, A. Polyak, T. Hayes, X. Yin, J. An, S. Zhang, Q. Hu, H. Yang, O. Ashual, O. Gafni, D. Parikh, S. Gupta, and Y. Taigman. Make-a-video: Text-to-video generation without text-video data. In The Eleventh International Conference on Learning Representations, 2023.

\section*{Genie：生成式交互环境}
S. Sudhakaran, M. González-Duque, C. Glanois, M. Freiberger, E. Najarro, and S. Risi. Promptguided level generation. In Proceedings of the Companion Conference on Genetic and Evolutionary Computation, pages 179–182, 2023.

A. Summerville, S. Snodgrass, M. Guzdial, C. Holmgård, A. K. Hoover, A. Isaksen, A. Nealen, and J. Togelius. Procedural content generation via machine learning (PCGML). IEEE Trans. Games, 10 (3):257–270, 2018.

G. Todd, S. Earle, M. U. Nasir, M. C. Green, and J. Togelius. Level generation through large language models. In Proceedings of the 18th International Conference on the Foundations of Digital Games, pages 1–8, 2023.

F. Torabi, G. Warnell, and P. Stone. Behavioral cloning from observation. arXiv preprint arXiv:1805.01954, 2018.

T. Unterthiner, S. van Steenkiste, K. Kurach, R. Marinier, M. Michalski, and S. Gelly. FVD: A new metric for video generation, 2019.

A. van den Oord, A. Razavi, B. Uria, Çağlar Ünlü, C. Nash, C. Wolff, C. Durkan, D. Ding, D. Górny, E. Gladchenko, F. Riedel, H. Qi, J. Kelly, J. Bauer, J. Donahue, J. Zhang, M. Malinowski, M. Bińkowski, P. Luc, R. Riachi, R. Strudel, T. P. I. Sander Dieleman, Y. Ganin, and Z. Eaton-Rosen. Imagen 2. URL https://deepmind.google/technologies/imagen-2/.

A. van den Oord, O. Vinyals, and K. Kavukcuoglu. Neural discrete representation learning. In Proceedings of the 31st International Conference on Neural Information Processing Systems, NIPS’17, page 6309–6318, Red Hook, NY, USA, 2017. Curran Associates Inc. ISBN 9781510860964.

A. Vaswani, N. Shazeer, N. Parmar, J. Uszkoreit, L. Jones, A. N. Gomez, L. Kaiser, and I. Polosukhin. Attention is all you need. In Advances in Neural Information Processing Systems, pages 5998–6008, 2017.

R. Villegas, M. Babaeizadeh, P.-J. Kindermans, H. Moraldo, H. Zhang, M. T. Saffar, S. Castro, J. Kunze, and D. Erhan. Phenaki: Variable length video generation from open domain textual descriptions. In International Conference on Learning Representations, 2023.

J. C. Walker, A. Razavi, and A. van den Oord. Predicting video with VQVAE, 2021.

Y. Wang, Y. He, Y. Li, K. Li, J. Yu, X. Ma, X. Chen, Y. Wang, P. Luo, Z. Liu, Y. Wang, L. Wang, and Y. Qiao. Internvid: A large-scale video-text dataset for multimodal understanding and generation, 2023.

L. Wong, G. Grand, A. K. Lew, N. D. Goodman, V. K. Mansinghka, J. Andreas, and J. B. Tenenbaum. From word models to world models: Translating from natural language to the probabilistic language of thought, 2023.

C. Wu, J. Liang, L. Ji, F. Yang, Y. Fang, D. Jiang, and N. Duan. Nüwa: Visual synthesis pre-training for neural visual world creation. In European conference on computer vision, pages 720–736. Springer, 2022.

M. Xu, W. Dai, C. Liu, X. Gao, W. Lin, G.-J. Qi, and H. Xiong. Spatial-temporal transformer networks for traffic flow forecasting. arXiv preprint arXiv:2001.02908, 2020.

W. Yan, Y. Zhang, P. Abbeel, and A. Srinivas. Videogpt: Video generation using vq-vae and transformers, 2021.

\section*{Genie：生成式交互环境}
W. Yan, D. Hafner, S. James, and P. Abbeel. Temporally consistent transformers for video generation. In A. Krause, E. Brunskill, K. Cho, B. Engelhardt, S. Sabato, and J. Scarlett, editors, Proceedings of the 40th International Conference on Machine Learning, volume 202 of Proceedings of Machine Learning Research, pages 39062–39098. PMLR, 23–29 Jul 2023.

M. Yang, Y. Du, K. Ghasemipour, J. Tompson, D. Schuurmans, and P. Abbeel. Learning interactive real-world simulators. arXiv preprint arXiv:2310.06114, 2023.

W. Ye, Y. Zhang, P. Abbeel, and Y. Gao. Become a proficient player with limited data through watching pure videos. In The Eleventh International Conference on Learning Representations, 2022.

L. Yu, Y. Cheng, K. Sohn, J. Lezama, H. Zhang, H. Chang, A. G. Hauptmann, M. Yang, Y. Hao, I. Essa, and L. Jiang. Magvit: Masked generative video transformer. In 2023 IEEE/CVF Conference on Computer Vision and Pattern Recognition (CVPR), pages 10459–10469, Los Alamitos, CA, USA, jun 2023. IEEE Computer Society. doi: 10.1109/CVPR52729.2023.01008.

\section*{Genie：生成式交互环境}
\section*{其它示例轨迹}

图 16 | 更多轨迹示例：模型可以接受手绘草图、文本到图像生成模型产生的图像，或真实照片作为提示；轨迹中的动作由人工输入提供。

\paragraph{动作提示} left / right / jump / no-op

图 17 | Platformers 中可控且一致的潜在动作：图中轨迹分别从 Platformers 数据集中的四个不同起始帧开始。每一列展示连续五次执行同一个潜在动作后的结果。尽管训练时没有动作标签，同一个潜在动作在不同提示帧上的效果仍较一致，并呈现出清晰语义：向左、向右、跳跃和无操作。

\section*{附录 B：数据集}
\subsection*{B.1 Platformers 数据集}

\paragraph{初始数据集}
作者从公开互联网视频中筛选出 2D 平台游戏相关视频，筛选条件如下：
\begin{itemize}
  \item 标题包含与 2D 平台游戏相关的关键词。
  \item 标题或描述中包含动作词，例如 \textit{speedrun} 或 \textit{playthrough}。
  \item 标题不包含 \textit{movie}、\textit{unboxing} 等否定关键词。
\end{itemize}

随后，作者将每个视频切分为 16 秒片段，并以 10 FPS 采样，因此每个片段包含 160 帧。得到的初始数据集约有 5500 万个视频片段，总时长约 24.4 万小时。关键词选择过程中，作者人工抽查检索结果，确认这些关键词通常能够返回 2D 平台游戏画面，而不是被大量共享相似关键词的其他视频淹没。

\paragraph{筛选流程}
作者注意到，初始数据集中有相当一部分视频质量较差，会影响模型性能。为此，论文采用类似 Baker et al. (2022) 的可扩展筛选方法，用学习得到的分类器进行数据清洗。高质量视频被定义为画面中有清晰游戏过程，且不包含菜单界面、主播头像等干扰元素。具体流程如下：
\begin{enumerate}
  \item 团队人工标注 1 万个视频，总耗时约 10 小时；标签从 5 分到 1 分，分别表示最好到最差。
  \item 训练一个 1100 万参数的 ResNet18 二分类器（He et al., 2016）。评分为 2 到 4 的样本被移除，只把 5 分样本视为好样本、1 分样本视为差样本。
  \item 根据模型预测结果及置信度设置保留规则，决定是否保留视频。
\end{enumerate}

这一结果与 Baker et al. (2022) 和 Oquab et al. (2023) 的观察一致：数据质量比单纯的数据数量更重要。尽管筛选后数据集的规模仅略高于原始数据集的 10\%，用筛选数据训练的模型在 FVD 上反而更好。最终数据集包含 680 万个视频片段，总时长超过 3 万小时。

表 4 | 数据筛选的影响。

\#Params FVD (↓)：原始数据集（5500 万视频）580M，61.4；筛选数据集（680 万视频）580M，54.8。

\section*{附录 C：训练细节}
\subsection*{C.1 潜在动作模型训练}

作者发现，增加码本大小，即增加可用动作数量，会带来一定收益；代价是人类和 AI 智能体试玩时的可操作性下降。

表 5 | Platformers 动作模型超参数。

\section*{参数值}
Encoder：num\_layers 20，d\_model 1024，num\_heads 16。Decoder：num\_layers 20，d\_model 1024，num\_heads 16。Codebook：num\_codes 8，patch\_size 16，latent\_dim 32。

模型输入被归一化到 0 到 1 之间，解码器的最终输出经过 sigmoid。

\subsection*{C.2 视频 tokenizer 训练}

本节说明视频 tokenizer 的训练设置。实验中，相比扩大编码器，扩大解码器更有效；增大批量大小也带来小幅收益，见表 6。

表 6 | Tokenizer 批量大小扩展超参数。

batch\_size / training hardware / FLOPs / PSNR：64 / 64 TPUv2 / $4.22 \times 10^{20}$ / 35.7；384 / 64 TPUv3 / $2.57 \times 10^{21}$ / 36.5。

表 7 | Platformers 视频 tokenizer 超参数。

\section*{参数值}
Encoder：num\_layers 12，d\_model 512，num\_heads 8，k/q\_size 64。Decoder：num\_layers 20，d\_model 1024，num\_heads 16，k/q\_size 64。Codebook：num\_codes 1024，patch\_size 4，latent\_dim 32。

视频 tokenizer 使用 AdamW 优化器训练 30 万步，并采用余弦衰减；优化器超参数见表 8。

表 8 | 视频 tokenizer 优化器超参数。

\section*{参数值}
max\_lr 3e-4，min\_lr 3e-4，$\beta_1$ 0.9，$\beta_2$ 0.9，weight\_decay 1e-4，warmup\_steps 10k。

\subsection*{C.3 动力学模型训练}

\section*{规模化实验细节}
本节补充规模化实验中的模型架构和计算预算。

\paragraph{模型规模扩展}
所有模型都使用批量大小 256，并训练 20 万步，因此每次运行约使用 750B 个训练 token。所有运行都采用批量并行和 stage-3 ZeRO 分片（Rajbhandari et al., 2020）；较大的模型还使用张量并行（Shoeybi et al., 2019）。该实验使用 TPUv2 和 TPUv3（Jouppi et al., 2020）。更多细节见表 10。

表 9 | 动力学模型优化器超参数。

\section*{参数值}
max\_lr 3e-5，min\_lr 3e-6，$\beta_1$ 0.9，$\beta_2$ 0.9，weight\_decay 1e-4，warmup\_steps 5k。

表 10 | 模型规模扩展中的架构和计算量。所有模型均使用批量大小 256 训练 20 万步，相当于 750B 个 token。

Parameters / num\_layers / num\_heads / d\_model / k/q size / training hardware / training time / FLOPs：41M / 18 / 8 / 512 / 64 / 64 TPUv2 / 3 days / $2.05 \times 10^{20}$；96M / 16 / 16 / 768 / 64 / 64 TPUv2 / 6 days / $3.58 \times 10^{20}$；192M / 20 / 18 / 1024 / 64 / 64 TPUv2 / 9 days / $6.4 \times 10^{20}$；404M / 21 / 12 / 1536 / 128 / 64 TPUv2 / 18 days / $1.2 \times 10^{21}$；811M / 20 / 20 / 2048 / 128 / 128 TPUv3 / 7 days / $2.2 \times 10^{21}$；1.6B / 28 / 22 / 2560 / 128 / 128 TPUv3 / 12 days / $4.04 \times 10^{21}$；2.7B / 36 / 22 / 3072 / 128 / 256 TPUv3 / 16 days / $6.91 \times 10^{21}$。

\paragraph{批量大小扩展}
所有模型都使用表 11 所示的 23 亿参数架构，并训练 20 万步。三次运行唯一的区别是硬件：批量大小 128、256、448 的模型分别在 64 个 TPUv3、128 个 TPUv3 和 64 个 TPUv5p 上训练。

表 11 | 批量大小扩展超参数。所有模型使用如下架构训练 20 万步，仅批量大小不同。

Parameters / num\_layers / num\_heads / d\_model / k/q size：2.3B / 34 / 20 / 2560 / 128。

\paragraph{Genie 模型}
最终 Genie 模型的动力学模型参数量、架构和计算量见表 12。作者训练了一个 101 亿参数的动力学模型，批量大小为 512，总共训练 12.5 万步，使用 256 个 TPUv5。

表 12 | Genie 动力学模型超参数。

Parameters / num\_layers / num\_heads / d\_model / k/q size / FLOPs：10.1B / 48 / 36 / 5120 / 128 / $6.6 \times 10^{22}$。

\section*{附录 E：行为克隆细节}

本节补充行为克隆实验的设置。作者在 Procgen CoinRun 环境（Cobbe et al., 2020）中训练，并在留出的测试集上评估。实验假设有一个由 R2D2 智能体（Kapturowski et al., 2018）生成的专家序列数据集，随后训练另一个智能体模仿这些序列。oracle 智能体可以访问对应的真实专家动作；下文说明如何用预训练 LAM 推断视频中的动作。

\subsection*{E.1 Genie LAM}

为让智能体从未见过的视频中学习模仿，可以使用在互联网视频上训练好的冻结 Genie LAM。给定专家序列 $\langle x_t, x_{t+1} \rangle$，模型提取对应的潜在动作标签 $a_t \leftarrow LAM(x_t, x_{t+1})$。随后训练策略 $\pi(a_t \mid x_t)$，预测专家在观测 $x_t$ 下采取潜在动作 $a_t$ 的概率。

这一流程与此前从视频中学习的工作类似（Baker et al., 2022; Torabi et al., 2018），但这些方法使用真实动作来标注视频，而这里使用完全离线学习得到的潜在动作。推断时，需要把策略输出的潜在动作映射到真实动作。为此，作者使用一小部分带动作标签的专家序列。给定专家序列 $\langle x_t, u_t, x_{t+1} \rangle$，其中 $u_t$ 表示真实动作，用于避免与预测潜在动作混淆，LAM 先得到潜在动作 $a_t$，再填充字典 $D$，把每个潜在动作映射到相应真实动作列表。于是，环境给出观测 $x_t$ 后，可以先从策略中取得最可能的潜在动作 $a_t \sim \pi(s_t)$，再从映射中采样真实动作 $u_t \sim D[a_t]$。

已有工作也曾从智能体策略产生的数据中学习潜在动作到真实动作的映射（Edwards et al., 2019; Ye et al., 2022）。本文发现，用专家数据建立映射更便于评估学得策略的质量。主文结果显示，智能体只需 200 个专家标签就能完成适配。

\subsection*{E.2 架构}

oracle 智能体和潜在 BC 智能体都使用 Transformer 作为策略。作者采用 ST-ViViT 编码帧序列 $\mathbf{x}_{1:t}=(x_1,\ldots,x_t)$。此前动作先经过 one-hot 表示，再与对应帧编码相加，作为加性嵌入。训练和推断时都使用长度为 4 的序列，批量大小为 16。

表 13 | BC 模型优化器超参数。

\section*{参数值}
max\_lr 3e-5，min\_lr 3e-6，$\beta_1$ 0.9，$\beta_2$ 0.96，weight\_decay 1e-4，warmup\_steps 5k。

表 14 | BC 策略超参数。

\section*{参数值}
Encoder：num\_layers 12，d\_model 512，patch\_size 4。Policy：linear\_layer 512。

oracle 和 Genie LAM 均用交叉熵损失训练，目标分别是真实动作和潜在动作。推断时，最终预测来自模型输出 logits 的采样。作者还发现，oracle 智能体在 10\% 的时间随机采样动作时表现更好。

\section*{附录 F：可复现实例}

本节给出一个自包含、可完整复现的案例研究。该设置可在单个中端 TPU/GPU 上训练，并在一周内完成。

\subsection*{F.1 数据收集}

首先需要收集模型训练数据。作者使用 Procgen 基准中的 CoinRun 环境（Cobbe et al., 2020），因为该环境包含数千个视觉上不同的关卡，动力学也接近简单的平台游戏。在 \textit{hard} 模式下，作者使用无动作重复的随机策略采集数据；关卡种子从 0 到 10000 中采样，每个关卡采集 1000 个时间步，总计 1000 万个转移。

\subsection*{F.2 视频 tokenizer 训练}

CoinRun 视频 tokenizer 沿用第 2.1 节的设置，并使用附录 C.2 中的优化器配置。该示例的主要区别是模型更小，见表 15；批量大小为 48 个序列，每个序列长度 16，因此每批共 768 张图像。这个配置可以放入单个 16G 内存的 TPU。模型在单个 TPU 上训练 3 天，可以完成 30 万步。

表 15 | CoinRun 视频 tokenizer 超参数。

\section*{参数值}
Encoder：num\_layers 8，d\_model 512，num\_heads 8。Decoder：num\_layers 8，d\_model 512，num\_heads 8。Codebook：num\_codes 1024，patch\_size 4，latent\_dim 32。

\subsection*{F.3 动力学模型与潜在动作模型训练}

视频 tokenizer 训练完成后，可以联合训练潜在动作模型和动力学模型。这里同样要求训练过程能够放入 16G 内存，因此使用批量大小 36，每个序列包含 16 帧，总计 576 张图像。潜在动作模型和动力学模型并行训练，设置沿用上文：潜在动作模型见附录 C.1，动力学模型见附录 C.3。二者使用表 9 中的优化器超参数并行训练 20 万步。实验中，该模型能够生成一致、可游玩的潜在动作，效果接近原始环境。

表 16 | CoinRun 动作模型超参数。

\section*{参数值}
Encoder：num\_layers 8，d\_model 512，num\_heads 8。Decoder：num\_layers 8，d\_model 512，num\_heads 8。Codebook：num\_codes 6，latent\_dim 32。

表 17 | CoinRun 动力学模型超参数。

\section*{参数值}
Architecture：num\_layers 12，d\_model 512，num\_layers 8。Sampling：temperature 1.0，maskgit\_steps 25。



\end{document}
